\originalchapter{8.06 作业6：自发辐射与光电离}
来源为Spring 2018 Assignment 6, Barton Zwiebach. 第1--3题有完整MIT题面, 中文转述并给\textbf{AI独立解答}. 第4、5题只有Griffiths 9.11、p.359及9.14、p.363题号, 仅保留引用, 不补造氢衰变教材题面.
\originalsection{第1题：三维谐振子的辐射寿命}
\begin{exercise}
带电粒子的各向同性三维谐振子第一激发态$\ket z, (\ket x\pm\ii\ket y)/\sqrt2$对应$\ell=1,m=0,\pm1$, 求向基态偶极自发辐射率和寿命. 粒子取电子, $\hbar\omega=3E_{\rm Ry}/4$求数值.
\end{exercise}
\begin{solution}
偶极耦合$H_I=-d\cdot E$, $d=q\boldsymbol r$. 体积$V$中每光子模式电场幅$\sqrt{\hbar\omega/(2\epsilon_0V)}$, 黄金规则乘模式密度$V\omega^2\dd\omega\dd\Omega/(2\pi)^3c^3$, 两横向偏振满足$\sum_\lambda|d\cdot\epsilon_\lambda|^2=|d|^2-|d\cdot\hat k|^2$. 角积分为$8\pi|d|^2/3$, 并用能量Dirac核的$1/\hbar$换频率, 得SI标准结果$\Gamma=\omega^3|d_{fi}|^2/(3\pi\epsilon_0\hbar c^3)$. 这是量子场单模输入与模式计数, 不把经典外加驱动说成可自行从真空辐射.

沿任一方向$\bra0r_i\ket i=\sqrt{\hbar/(2m\omega)}$, 其余0. 三个角动量组合的矢量模平方都$|d|^2=q^2\hbar/(2m\omega)$, 所以$\Gamma=q^2\omega^2/(6\pi\epsilon_0mc^3)$,$\tau=6\pi\epsilon_0mc^3/(q^2\omega^2)$, 旋转对称使$m$寿命相同. 题给$E_{\rm Ry}=13.6$eV, $\hbar\omega=10.2$eV, 得$\omega\simeq1.54965\times10^{16}$s$^{-1}$, $\Gamma\simeq1.50482\times10^9$ s$^{-1}$,$\tau\simeq6.64532\times10^{-10}$s. 能差与氢2P--1S相同, 但此处矩阵元属于谐振子, 寿命不能直接当成氢2P寿命.
\end{solution}
\originalsection{第2题：一维深井的光电离}
\begin{exercise}
方井$0<x<a$势$-V_0$、其余0且很深, 基态近似$\sqrt{2/a}\sin(\pi x/a)$、能量近似$-V_0$. 外场$E(t)=2E_0\cos\omega t$沿$x$. 末态近似长度$L$盒中平面波, (a)求一维末态密度及近似适用条件. (b)算总电离率, 检查$L$消去. 原件最后势支写$x>0$与井内重叠, 按一致方井定义更正为$x>a$.
\end{exercise}
\begin{solution}
连续末能$E_f=\hbar\omega-V_0>0$, $p=\sqrt{2mE_f}$. 每一传播方向的态密度$\rho_E^{(\pm)}=Lm/(2\pi\hbar p)$, 两方向相加$Lm/(\pi\hbar p)$. 吸收阈值$\hbar\omega>V_0$只使末态存在; 用弱散射平面波还要求$E_f\gg V_0$, 即$\hbar\omega\gg2V_0$. 初态近无限井须$V_0\gg\hbar^2/(ma^2)$.

电子扰动$eE(t)x$, 吸收Fourier分量$eE_0x$, 不能把完整$2E_0$再当吸收幅而多乘4. 令$k=p/\hbar,b=\pi/a$, $J(k)=\int_0^a x\sin(bx)\e^{-\ii kx}\dd x$, 矩阵元$eE_0\sqrt{2/(aL)}J$. $F(k)=\int_0^a\sin(bx)\e^{-\ii kx}\dd x=b(1+\e^{-\ii ka})/(b^2-k^2)$, $J=\ii F'$, 所以
\[
J=b\left[\frac{a\e^{-\ii ka}}{b^2-k^2}
+\frac{2\ii k(1+\e^{-\ii ka})}{(b^2-k^2)^2}\right],\qquad
\Gamma=\frac{4e^2E_0^2m}{a\hbar^2p}|J(k)|^2.
\]
两方向有相同模, 所列$\rho_E$已经计两向, 不重复加倍. $L$从矩阵元和密度抵消. $k=b$闭式表面极点是可去的, 应用积分或极限评价; 实近似高能本来不在该点. 自由末态与真实井束缚态不严格正交, 故低能精确偶极和原点平移性质需要完整散射态, 不能将高能近似推到阈值.
\end{solution}
\originalsection{第3题：热光子的两种发射}
\begin{exercise}
求温度$T$热场的自发/受激发射比, 从光子Boltzmann分布推平均占据数. 室温何频率两者相同? 可见光与maser哪个过程主要?
\end{exercise}
\begin{solution}
单模$n$光子概率$P_n=(1-z)z^n$,$z=\e^{-\hbar\omega/(k_BT)}$, 几何级数及参数导数给$\bar n=z/(1-z)=1/(\e^{\hbar\omega/(k_BT)}-1)$. 发射的创建算符幅$\sqrt{n+1}$, 率分成$n$受激与1自发项, 比$r=1/\bar n=\e^{\hbar\omega/(k_BT)}-1$. 两者相同时$\nu=k_BT\log2/h$, 300K时约$4.333\times10^{12}$Hz（远红外）. 可见光例如$5\times10^{14}$Hz有指数极大$r$, 热场受激很小; 微波10GHz时$r\simeq0.00160$, 受激占优. 这里比较\textbf{热平衡背景}中的同一跃迁, 激光强外加非热光场的受激项可另大幅增强.
\end{solution}
